\section{A strong baseline, not a straw man}
\label{sec:baseline}

The Lean reference describes \grind{} as cooperating proof-producing engines
organised around contradiction, congruence closure, constraint propagation,
E-matching, guided case analysis, and theory solvers. That is already a
shared-fact architecture, not a sequence of independent tactic
calls~\cite{grind-ref,grind-paper}. A credible successor must reuse those
strengths rather than rediscover them.

The baseline must also be \emph{dated}. Lean v4.34.0 was released on
14~September 2026~\cite{lean-release-434}. Designing against a remembered
version of \grind{} is the most common way to overstate a contribution.
Table~\ref{tab:baseline} records capabilities that the pinned configuration
and implementation sources show are already present, and what each implies for
this design.

\begin{table}[htbp]
\centering\small
\begin{tabularx}{\textwidth}{@{}p{0.17\textwidth}Xp{0.30\textwidth}@{}}
\toprule
Area & Already present in the inspected source & Consequence for this design \\
\midrule
Equality &
Congruence reasoning, function-valued congruence, extensionality controls,
injectivity support &
Do not advertise ordinary e-graphs or function equalities as new. \\
Arithmetic &
Ring and field algebra, linear arithmetic, integer arithmetic,
associative/commutative reasoning, order support, model-based theory
combination &
Add certified inequality reasoning and structure, not a duplicate ring
normaliser. \\
Quantifiers, search &
E-matching, term-generation limits, case splits, lookahead &
Distinguish goal-directed \emph{construction} from existing instantiation and
splitting. \\
Bit-vectors &
\code{grind => bv\_decide} encodes relevant equivalence classes and facts
learned by theory solvers and E-matching into the SAT problem &
Basic cooperation between bit-blasting and \grind{} is existing
functionality, not a \forge{} innovation. \\
Conditional rewriting &
\grind{} discharges side conditions for conditional \code{Sym.simp} theorems &
The proposal is broader cross-worker obligation scheduling, not the invention
of guarded rewriting. \\
Libraries &
Configurable library suggestions and local-definition use &
Improve retrieval and obligation-driven use; do not claim retrieval is wholly
absent. \\
Engineering &
Saved states, proof-related diagnostics, explicit search limits &
Reuse existing machinery behind a versioned adapter. \\
\bottomrule
\end{tabularx}
\caption{Baseline features from the pinned Lean v4.34.0 configuration and
implementation sources~\cite{grind-config,grind-types,grind-lia,grind-ematch,lean-release-434}.
The table is not a completeness claim for any theory.}
\label{tab:baseline}
\end{table}

For reproducible comparison, the pinned configuration source records defaults
including nine search-branch splits, five E-matching rounds before a split,
generation bounds of eight, and one thousand E-matching instances per
branch~\cite{grind-config}. Many further limits are exposed. A future
benchmark should record the \emph{effective options}, not merely write
``default \grind{}'', and should include a reasonably tuned baseline as well
as the stock one.

\subsection{The neighbouring tactics, and the gap between them}
\label{sec:tactic-gap}

\grind{} is the baseline this design builds on, but it is not the only tactic a
new capability has to be measured against. Four others bound the space, and
naming them is what makes the remaining gap precise rather than rhetorical.

\begin{longtable}{@{}>{\raggedright\arraybackslash}p{0.17\textwidth}>{\raggedright\arraybackslash}p{0.38\textwidth}>{\raggedright\arraybackslash}p{0.38\textwidth}@{}}
\toprule
Tactic & What it does & What it does not do \\
\midrule
\endhead
\code{exact?} / \code{apply?} &
Finds an \emph{existing} library lemma that closes the goal. &
Cannot \emph{compose} a new term. Fails on
$(A\to B\to C)\to(A\to B)\to A\to C$ unless that exact lemma exists. \\
\code{tauto} / \code{itauto} &
Decides propositional goals; \code{itauto} is intuitionistically complete. &
\code{Prop} only; needs Mathlib imported; no negative verdicts surfaced as
such; nothing for \code{Type}. \\
Aesop &
General rule-based proof search. &
Heuristic; no non-inhabitation answers.~\cite{aesop} \\
\code{decide} &
Decidable ground propositions. &
Nothing polymorphic or data-level. \\
\code{omega} &
Linear integer and natural arithmetic. &
The reference explicitly does not describe it as complete for every
integer-linear problem.~\cite{tactic-reference} \\
\bottomrule
\end{longtable}

The gap this leaves is \emph{higher-order plumbing that has to be constructed
rather than retrieved}: currying and uncurrying, projections, composition,
distribution laws such as $A\times(B\oplus C)\to(A\times B)\oplus(A\times C)$,
and continuation shuffles such as $((A\to B)\to A)\to(A\to A)$. These are terms
one writes constantly. \code{exact?} cannot build them, \code{tauto} covers only
the \code{Prop} shadow of them, and Aesop will search for them without ever
being able to report that a particular one does not exist.

That last point is the one worth dwelling on, because it is a capability and
not merely a diagnostic. \S\ref{sec:negative} argues that a tactic which can
sometimes say \emph{provably no such term exists} is doing something none of
the five above does.

\subsection{Related systems: reuse, do not relabel}

Aesop already provides rule-based AND/OR proof search, distinguishes
normalisation from safe and unsafe rules, and uses best-first exploration.
Its documented induction examples still perform induction explicitly. It is a
source of search infrastructure, not a feature to rediscover~\cite{aesop}.

Duper supplies proof-producing saturation over supplied facts, with manual
premise selection and a configurable portfolio. It is a complementary
clause-search worker, especially after a relevance filter has selected a small
premise set~\cite{duper}.

Mathlib's \code{linarith} separates certificate search from proof
reconstruction, and \code{nlinarith} adds specific nonlinear preprocessing.
That is the right architectural precedent, but it is not a general
sum-of-squares solver~\cite{mathlib-linarith}. Foundational checking of
numerically discovered SOS certificates has a substantial history; Harrison's
HOL Light work explicitly separates sophisticated search from relatively
straightforward verification~\cite{harrison-sos}.

There is already a Lean \code{sos} project implementing proof-producing
nonlinear reasoning with rational certificates. Its documented architecture
separates reification, search, checking, and verification; named interfaces
include \code{SOS.Engine.solve} and \code{SOS.Engine.Result.check}, and
\code{sos?} can emit a replayable witness. Its computational polynomial
substrate is \code{Hex.MvPoly}~\cite{lean-sos,sos-engine,sos-core,hex-mv-poly}.
These are concrete integration targets. The \forge{} contribution is
\emph{when} to request such a certificate and \emph{how} to return its proved
consequences to structural and saturation workers --- not the invention of
sum-of-squares certificates.

The historical \code{lean-egg} project should not become a mandatory new
dependency: its repository is marked as no longer developed and directs users
towards \grind{}~\cite{lean-egg}. The useful idea is bounded, proof-producing
equality search, not allegiance to that
package~\cite{egg,egg-lean-paper,eqsat-paper}.

\begin{remark}[A version caveat that matters]
The nine contributing efforts observed the Lean \code{sos} project's
\texttt{lean-toolchain} pinning \code{v4.32.0-rc1}, against a \code{v4.34.0}
baseline~\cite{sos-toolchain}. These projects do not automatically form a
build-compatible package set on the day a new Lean release appears. Pin the
core toolchain, select mutually compatible Mathlib and optional-backend
revisions, and keep their adapters separate. No \code{lake-manifest.json} is
supplied here for a combination that was never built.
\end{remark}

\subsection{A dependency map with integration decisions}

\begin{longtable}{@{}>{\raggedright\arraybackslash}p{0.17\textwidth}>{\raggedright\arraybackslash}p{0.38\textwidth}>{\raggedright\arraybackslash}p{0.38\textwidth}@{}}
\toprule
Component & What it provides & Integration decision \\
\midrule
\endhead
Lean core / \grind{} &
Shared equality and local-theory reasoning, proof terms, saved states. &
Keep as the default local closer; start at the tactic boundary, move to a
version-pinned adapter only where measurement justifies
it.~\cite{grind-ref,grind-main} \\
Aesop &
Best-first AND/OR proof search with indexed rules. &
Reuse rule-search ideas and adapters; do not claim AND/OR search is
novel.~\cite{aesop} \\
Mathlib tactics &
\code{simp}, \code{ring}, \code{linarith}, \code{nlinarith},
\code{positivity}, \code{linear\_combination}, cast normalisation. &
Cheap proof reconstruction and preprocessing before expensive
search.~\cite{mathlib-ring,mathlib-linarith,mathlib-positivity,mathlib-linear-combination} \\
Lean \code{sos} &
Rational nonlinear certificates and a proof-producing tactic. &
The production nonlinear backend. Do not replace it with the prototype's
restricted LP dictionary.~\cite{lean-sos} \\
Lean \code{lp} &
Verified linear programming. &
Candidate backend for the linear feasibility problems underlying cone and
witness search.~\cite{lean-lp,lp-blog} \\
Duper &
Proof-producing saturation inside Lean. &
An optional, separately pinned first-order worker over a small retrieved
premise set --- never an unchecked external verdict.~\cite{duper} \\
\code{lean-auto} &
Translation and monomorphization infrastructure. &
Reuse rather than invent an unverified dependent-to-first-order
erasure.~\cite{lean-auto} \\
\code{lean-smt} / cvc5 &
SMT search with Lean reconstruction; unsupported steps remain goals. &
Optional backend with an explicit residual-obligation
interface.~\cite{lean-smt,cvc5-proofs} \\
QuerySMT / LeanHammer &
Hint-based reconstruction and premise selection. &
Prior art for the retrieval and reconstruction lanes.~\cite{query-smt,lean-hammer} \\
PBLean &
Pseudo-Boolean proof certificates for Lean~4. &
Reference point for the Boolean certificate lane.~\cite{pblean} \\
Z3 &
Candidate models and quantified search techniques. &
Untrusted search assistance; a verdict is not a Lean
proof.~\cite{z3-quantifiers} \\
\code{egg} / CCLemma &
Equality saturation and inductive-lemma-search prior art. &
Algorithmic references and possible isolated search
processes.~\cite{egg,cclemma} \\
Canonical &
Dependent type-theory inhabitation and unification. &
Integrate by its proof-term interface, not by flattening away the
dependencies that make its task meaningful.~\cite{canonical,canonical-min} \\
\bottomrule
\end{longtable}

\subsection{A portfolio alone is not the proposal}

Running \code{first | grind | aesop | nlinarith} is a portfolio. The important
difference is that in \forge{} an \emph{unsuccessful} attempt still produces
structured, reusable information: which recursive argument changed, which
equality would enable progress, which witness template was insufficient, and
which arithmetic side condition remains. That feedback directs the next
attempt, and every checked intermediate result returns to the shared proof
state.

A benchmark must therefore include a same-workers, same-budget,
\emph{non-communicating} portfolio as a baseline. That is the experiment that
distinguishes ``integration helps'' from ``we added more solvers''
(Section~\ref{sec:evaluation}).
